\subsection{Vínculo explícito de certificado del Owner}\label{sec:impl-vinculo-certificado}

El vínculo relaciona una cuenta Owner consigo misma con un despliegue, una raíz, una revisión de confianza y un certificado público de la CA interna (\cref{fig:cripto-pki}). Sobre él se apoya el primer factor opcional por certificado de CU-03 (\cref{tab:cu-03}). El vínculo y el acceso usan declaraciones firmadas de longitud fija con propósitos separados (\cref{tab:anexo-canones}): \texttt{OWNCERT1}, de 150 bytes, registra o retira el vínculo con revisión de cuenta y generación; por sí solos, sus valores no autentican ni acreditan una escritura durable.

\paragraph{Verificador.} Bajo la sesión MFA y antes de los bloqueos, el perfil Partner exige RSA de 3072 bits, exponente 65537, SHA-256/RSA y usos de clave exactos con clientAuth y emailProtection, que el perfil de la declaración de participantes rechaza (\cref{sec:impl-declaracion-participante,sec:pki-declaraciones}). Recalcula el resumen SHA-256, verifica firma, raíz y CRL obligatoria al instante pedido y compara confianza y huellas: prueba consistencia con la confianza aportada, no su publicación ni su vigencia futura.

\paragraph{Persistencia y consulta.} PostgreSQL toma el bloqueo compartido de auditoría y el de la cuenta, relee Owner activo, contadores y confianza publicada y revalida la vigencia con el reloj leído tras las esperas. Alta y retiro solo añaden filas, sin modificar las anteriores, con su evento en la misma transacción; los disparadores no verifican RSA ni autentican. Cada cuenta tiene como máximo un vínculo sin retirar y una huella no pasa a otro Owner. El retiro es definitivo, conserva la captura y procede aunque el certificado haya vencido. La consulta del vínculo actual devuelve el recibo o indica que no hay vínculo, y falla si encuentra más de uno sin retirar; no audita ni verifica RSA o confianza vigente, así que su respuesta no impide un retiro posterior ni acredita vigencia.

\paragraph{Vista \emph{Mi certificado}.} Usa un UUID por intención y un único PEM público; firma con el ayudante local, sin que clave ni contraseña salgan del navegador, o importa una firma externa (\cref{sec:impl-criptografia-local}). Una escritura incierta se captura antes de esperar y se concilia por UUID exacto, sin repetirla salvo decisión expresa (\cref{sec:impl-qadra-conciliacion}); al reingresar solo queda en memoria el material público admitido (\cref{sec:impl-qadra-borradores}), y el recibo descargable no presenta la confianza histórica como vigente.

\paragraph{Acceso por certificado con MFA.} \texttt{OWNAUTH1} (\cref{tab:anexo-canones}) fija cuenta Owner, generación, vínculo y confianza; su intervalo de 1 a 300 s, con vencimiento exclusivo, solo limita presentar la prueba, y la validez de certificado y confianza es otra frontera. La captura en Redis vence en el instante firmado (\cref{fig:impl-acceso-certificado}). La revisión de cuenta puede avanzar al consumir un código de recuperación, pero generación y procedencia deben coincidir. La sesión guarda la revisión de confianza y sus plazos quedan acotados por la vigencia de la credencial; los registros de contraseña no cambian (ADR 0068).

\begin{figure}[!htb]
  \centering
  \begin{tikzpicture}[yscale=0.79,
  tramo/.style={draw=ESCOMGris, dotted, line width=0.5pt},
  tramoL/.style={font=\diag\bfseries, text=ESCOMAzulOscuro, fill=white,
    anchor=west, inner sep=1.5pt}]
  \Actores{usr/1.5/{Owner\\ (Qadra)},api/5.6/{API},rds/9.9/{Redis},pg/14.3/{PostgreSQL}}{-12.55}

  % ==== Declaración ====
  \draw[tramo] (0,-0.85) -- (16,-0.85);
  \node[tramoL] at (2.95,-0.85) {Declaración nueva};
  \Msg{1.5}{5.6}{-1.35}{\textbf{1} ¿acceso habilitado?}
  \Msg[msgR]{5.6}{1.5}{-1.8}{solo un booleano}
  \Paso{1.9}{-1.88}{3.5}{Elige el recibo público de su vínculo; conserva cuenta, vínculo, huella y nombre}
  \Msg{1.5}{5.6}{-3.3}{\textbf{2} pide declaración}
  \Msg{5.6}{9.9}{-3.75}{cuotas de inicio}
  \draw[msg] (5.6,-4.2) -- node[msgL, pos=0.74]{Owner activo, vínculo y confianza} (14.3,-4.2);
  \Msg{5.6}{9.9}{-4.65}{captura de un uso; vence al instante firmado}
  \Msg[msgR]{5.6}{1.5}{-5.1}{\texttt{OWNAUTH1} (182 B)}
  \Paso{1.9}{-5.25}{3.5}{Firma los bytes exactos con el ayudante local u OpenSSL}

  % ==== Prueba ====
  \draw[tramo] (0,-6.5) -- (16,-6.5);
  \node[tramoL] at (2.95,-6.5) {Prueba y segundo factor};
  \Msg{1.5}{5.6}{-7.0}{\textbf{3} firma de 384 B}
  \Msg{5.6}{9.9}{-7.45}{cuotas de prueba; consume la captura}
  \Paso{5.6}{-7.6}{3.6}{Verifica RSA (perfil Partner), confianza e instante}
  \Paso[pasoError]{9.9}{-7.6}{3.8}{Firma alterada o repetida: rechazo; la captura ya se consumió}
  \draw[msg] (5.6,-8.95) -- node[msgL, pos=0.74]{revalida cuenta, vínculo y confianza} (14.3,-8.95);
  \Msg{5.6}{9.9}{-9.4}{desafío MFA}
  \Msg[msgR]{5.6}{1.5}{-9.85}{desafío, nunca una sesión}
  \Msg{1.5}{5.6}{-10.4}{\textbf{4} código TOTP o de recuperación}
  \Msg{5.6}{9.9}{-10.85}{sesión con procedencia del vínculo}
  \Msg[msgR]{5.6}{1.5}{-11.3}{token de sesión}
  \Paso{1.9}{-11.45}{3.5}{Qadra exige el mismo Owner elegido y con rol Owner}
  \Paso{9.9}{-11.0}{3.8}{Retirar el vínculo o cambiar la confianza la invalida, sin pasarla a contraseña}
\end{tikzpicture}

  \caption{Acceso del Owner por certificado. La firma reemplaza el primer factor y solo concede el desafío MFA del acceso por contraseña; cuenta, vínculo y confianza se revalidan antes del desafío y de la sesión.}
  \label{fig:impl-acceso-certificado}
\end{figure}

\paragraph{Servicio y habilitación.} Las rutas del vínculo y del acceso solo admiten JSON estricto y acotado con evidencia pública, comparten admisión y presupuesto con la API y no reintentan escrituras inciertas; la preparación entrega los bytes exactos que se firman, y el servicio no retiene el token ni modifica campos de usuario. El binario no ejecuta migraciones ni publica confianza, y el alta exige la confianza publicada vigente; registrar un vínculo no enrola cuentas ni habilita el acceso por certificado ni la firma individual, que usa ese vínculo (\cref{sec:impl-firma-individual}). El acceso por certificado es opcional y está desactivado por omisión; sus cuotas y tiempos se validan al arrancar (límites en la \cref{sec:pruebas-vinculo-certificado}) y, si la configuración está incompleta, sus rutas no existen y una ruta pública anuncia que no está disponible. El navegador no puede habilitarlo, y habilitarlo no registra titulares, no publica confianza ni entrega claves.

\paragraph{Interfaz.} \emph{Ingresar con certificado} conserva la contraseña si la disponibilidad es negativa, ausente o inválida, y no usa el recibo elegido como prueba de vigencia. Cada intento firma una declaración nueva, uno incierto o vencido exige otra, y el plazo local monotónico no sustituye al del servidor (\cref{sec:impl-qadra-sesion}).

\paragraph{Límites.} Tras restaurar una copia, una firma nueva solo concede MFA nueva. La admisión de inventario no sustituye el verificador global de auditoría ni acredita una restauración completa. La habilitación privada está en la \cref{sec:impl-portal-ca}, y la verificación, con firmas OpenSSL y no del ayudante Rust/WebAssembly, en la \cref{sec:pruebas-vinculo-certificado}.

\Fuentes{ADR 0067 y 0068; \rutarepo{docs/owner-certificate-login-http.md}; \rutarepo{docs/owner-certificate-login-operations.md}; \rutarepo{docs/owner-certificate-login-interface.md}.}
